\scanpage{103}
Mais il faut encore expliciter la structure de $\unsure{S|\mathcal S}\Pi M_G$, en termes des $M_{g_\alpha,\widehat I_\alpha}$ — nous le ferons au n° suivant.

Pour un morphisme de normalisation quelconque (pas néc strict), le diagramme commutatif donne un diagramme comm (à iso. de foncteurs près)
\[
(146)\qquad
\diagram{\begin{tikzcd}[ampersand replacement=\&, column sep=small,row sep=large]
\Pi_1M_{\deleted{\unsure{G'}}G} \arrow[rr,leftrightarrow] \arrow[d,"\unsure{\sim|\wr}"] \&\& \Pi_1M_{G'} \arrow[d,"\unsure{\sim|\wr}"] \\
\Pi_1M_{\widetilde G}=\prod_{\alpha\in S}M_{g_\alpha,\widehat I_\alpha}
\arrow[rr] \arrow[dr,"{\substack{\text{on identifie }S'\\\text{à une partie de }S,\\\text{et }g_\alpha\text{ à }g'_{\alpha'}}}"'] \&\&
\Pi_1M_{\widetilde G'}=\prod_{\alpha'\in S'}\deleted{\unsure{1}}M_{g'_{\alpha'},\widehat I'_{\alpha'}} \\
\& \prod_{\alpha'\in S'\deleted{\unsure{*}}}M_{g_\alpha,\widehat I_\alpha} \arrow[ur,"\psi"'] \&
\end{tikzcd}}
\]
\margin{Ceci provient d'un diagramme analogue au niveau des multiplicités modulaires elles-mêmes\ldots}

La flèche oblique $\psi$ est induite par des flèches
\[
(147)\qquad \Pi_1M_{g_\alpha,\widehat I_\alpha}\longrightarrow\Pi_1M_{g_\alpha^{\deleted{\unsure{'}}},\widehat I'_\alpha}
\]
compte tenu du fait qu'on a $\widehat I'_\alpha\subset\widehat I_\alpha$.

Ainsi la connaissance des morphismes de normalisation \inserted{\unsure{généraux|général}} est \unsure{réduite|réduit} à celle des morphismes
\[
(148)\qquad \Pi_1M_{g,I}\longrightarrow\Pi_1M_{g,I'}\quad\text{pour }I'\subset I,
\]
\margin{102}
